\subsection{Country of Fools}

\subsubsection{Описание решения}

По условию задачи необходимо вывести последовательность чисел, в которой будет больше всего пар различных элементов, идущих подряд. Если выводить числа циклично \textit{(а-ля round-robin)}, то в конце мы можем получить довольно длинную последовательность одинаковых элементов.

Исправим стратегию. Для начала отсортируем массив по убыванию количества знаков каждого типа. Заметим, что первые элементы массива при желании можно разделить на две категории:
\begin{itemize}
  \item \textit{знак-вышка (pillar)} --- обладает уникальным значением количества, визуально выглядит как одиночная вышка;
  \item \textit{знаки-равнины (plains)} --- имеют общее для группы значение количества, визуально похожи на длинную \( (\geq 2) \) ступеньку. 
\end{itemize}

В алгоритме \textit{round-robin} на последнем этапе мы могли получить длинные цепочки элементов знака-вышки. Чтобы потенциально её уменьшить, нужно чередовать на каждой итерации вывод знака-вышки со знаками уровня ниже до тех пор, пока он не присоединятся к знакам-равнинам.

Если на данной итерации знак-вышка слилась со знаками-равнинами, то на следующих итерациях достаточно вернутся к циклическому алгоритму: новых знаков-вышек не образуется. Отметим, что если на первой итерации первые элементы --- знаки-равнины, то алгоритм \textit{round-robin} отработает корректно.

Для реализации описанного алгоритма выполним предобработку отсортированного массива, чтобы сгруппировать знаки с одинаковым значением количества. Для удобства полученный артефакт будем хранить в двусвязном списке из двусвязных списков. Если какой-либо из списков пустеет, удаляем его из главного списка. Алгоритм завершится, когда главный список опустеет.

\subsubsection{Асимптотическая сложность}

\textit{Временная сложность решения} --- \( \mathcal{O}(K\cdot\log K+N) \). Сортировка исходного массива из \( K \) элементов занимает линейно-логарифмическое время. Последующая обработка массива по заданному алгоритму --- линейная, поскольку в сумме для вывода \( N \) элементов используется константное число операций, каждая из которых занимает константное время в случае двусвязного списка.

\textit{Пространственная сложность решения} --- \( \mathcal{O}(K) \). В памяти хранятся два двусвязных списка из \( K \) элементов каждый. Числа выводятся в потоковом режиме, поэтому требуют константное количество памяти.

\subsubsection{Листинг кода}

\begin{lstlisting}
#include <cstddef>
#include <forward_list>
#include <functional>
#include <iostream>
#include <iterator>
#include <list>
#include <stdexcept>

class Sign {
public:
  std::size_t count = 0;
  std::size_t type = 0;

  // use this instead of `bool operator>(const Sign& other) const { ... }`
  // otherwise i would have to define getters (POD vs Non-POD)
  friend bool operator>(const Sign& lsn, const Sign& rsn) {
    return lsn.count > rsn.count;
  }
};

namespace {
std::list<std::list<Sign>> GenerateList(const std::forward_list<Sign>& signs) {
  if (signs.begin() == signs.end()) {
    throw std::invalid_argument("forward list is empty");
  }
  std::list<std::list<Sign>> list = {{*signs.begin()}};
  for (auto i = ++signs.begin(); i != signs.end() && i->count != 0; i++) {
    if (list.back().front().count == i->count) {
      list.back().push_front(*i);
      continue;
    }
    list.push_back({*i});
  }
  return list;
}

void StepPillar(std::list<std::list<Sign>>& list) {
  // taking item from a pillar
  auto next_list = std::next(list.begin());
  auto top = list.begin()->begin();
  std::cout << top->type << " ";
  top->count--;
  // if pillar is the only one item
  if (next_list == list.end()) {
    if (top->count == 0) {
      list.pop_front();
    }
    return;
  }
  // if pillar has transformed into plain, extend next item
  if (next_list->front().count == top->count) {
    next_list->push_front(*top);
    list.pop_front();
  }
  // take item from next plain
  std::cout << next_list->back().type << " ";
  next_list->back().count--;
  const Sign hanging_sign = next_list->back();
  next_list->pop_back();
  // if hanging sign was the pillar
  if (next_list->empty()) {
    next_list = list.erase(next_list);
  } else {
    ++next_list;
  }
  if (hanging_sign.count == 0) {
    return;
  }
  // extend next plain, if possible
  if (next_list != list.end() && hanging_sign.count == next_list->front().count) {
    next_list->push_front(hanging_sign);
    // otherwise, create a new pillar
  } else {
    list.insert(next_list, {hanging_sign});
  }
}

void ProcessPlain(std::list<std::list<Sign>>& list) {
  auto next_list = std::next(list.begin());
  auto plain = list.begin();
  while (!plain->empty()) {
    std::cout << plain->back().type << " ";
    plain->back().count--;
    // if item is not gone yet
    if (plain->back().count > 0) {
      // extend next plain, if possible
      if (next_list != list.end() && plain->back().count == next_list->front().count) {
        next_list->push_front(plain->back());
        // otherwise, create a new pillar
      } else {
        next_list = list.insert(next_list, {plain->back()});
      }
    }
    plain->pop_back();
  }
  list.pop_front();
}
}  // namespace

int main() {
  std::size_t k_types = 1;
  std::cin >> k_types;
  std::forward_list<Sign> signs;
  for (std::size_t i = 0; i < k_types; i++) {
    Sign sign = {0, i + 1};
    std::cin >> sign.count;
    signs.push_front(sign);
  }
  signs.sort(std::greater<>());

  std::list<std::list<Sign>> list = {};
  try {
    list = GenerateList(signs);
  } catch (const std::invalid_argument& e) {
    std::cout << "Problem constraints are not met.";
    return 1;
  }
  signs.clear();

  while (!list.empty() && std::next(list.front().begin()) == list.front().end()) {
    StepPillar(list);
  }
  while (!list.empty()) {
    ProcessPlain(list);
  }
  return 0;
}
\end{lstlisting}